% !TeX root = ../main.tex
\section{Расчет нелинейных характеристик усилителя по постоянному току}

Для исследования поведения каскада при больших уровнях сигнала и корректного выбора режима по постоянному току используется нелинейная SPICE-модель транзистора NE68133 (\texttt{NE68133Package.lib}), учитывающая паразитные параметры корпуса (индуктивности выводов и межэлектродные емкости), а также нелинейные свойства переходов.

Принципиальная схема каскада для снятия статических вольт-амперных характеристик (ВАХ) приведена на рис.~\ref{fig:dc_schematic}.

\begin{figure}[H]
    \centering
    \includegraphics[width=0.75\textwidth]{02_sch.png}
    \caption{Схема каскада на транзисторе NE68133 с питанием по постоянному току}
    \label{fig:dc_schematic}
\end{figure}

Подача постоянных смещений осуществляется через радиочастотные развязывающие дроссели $L_4 = L_5 = 10\text{ nH}$, предотвращающие утечку высокочастотного сигнала в цепи источников питания. Разделительные конденсаторы $C_2 = C_3 = 10\text{ pF}$ устраняют протекание постоянной составляющей тока через внешние 50-омные нагрузки.

В среде Ansys Designer был выполнен расчет по постоянному току (\textit{DC Analysis}) со следующими параметрами развертки:
\begin{itemize}
    \item напряжение базы $V_B$: диапазон $0{,}75\dots 0{,}85\text{ В}$ с шагом $0{,}01\text{ В}$;
    \item напряжение коллектора $V_C$: диапазон $0\dots 5{,}0\text{ В}$ с шагом $0{,}02\text{ В}$.
\end{itemize}

Результаты моделирования выходных вольт-амперных характеристик $I_C(V_C)$ представлены на рис.~\ref{fig:dc_iv}.

\begin{figure}[H]
    \centering
    \includegraphics[width=0.78\textwidth]{02_plot.png}
    \caption{Семейство выходных вольт-амперных характеристик транзистора NE68133 при вариации напряжения базы $V_B$}
    \label{fig:dc_iv}
\end{figure}

Анализ полученных характеристик показывает:
\begin{itemize}
    \item Граничное напряжение насыщения транзистора составляет $V_{C,\text{нас}} \approx 0{,}25\dots 0{,}3\text{ В}$.
    \item В активной области характеристики имеют выраженный линейный наклон, обусловленный эффектом Эрли (модуляцией ширины базы).
    \item Заданный в исходных данных малосигнальный режим $V_{CE} = 2{,}5\text{ В}$, $I_C = 3\text{ мА}$ надежно лежит в середине пологой активной усилительной зоны (соответствует напряжению базы $V_B \approx 0{,}78\dots 0{,}79\text{ В}$), что гарантирует высокую линейность усиления и отсутствие насыщения при малых входных уровнях.
\end{itemize}
